\section{Memoria Virtual}

La \textit{memoria virtual} es una técnica de administración de memoria que permite utilizar el almacenamiento secundario como una extensión de la memoria principal. De esta forma, un proceso puede disponer de un espacio de direcciones mayor que la memoria RAM físicamente instalada.

Para lograrlo, se distingue entre las direcciones generadas por el programa (\emph{direcciones virtuales}) y las direcciones reales de la memoria (\emph{direcciones físicas}). Cada vez que un proceso accede a memoria, el hardware y el sistema operativo traducen automáticamente la dirección virtual a la dirección física correspondiente.

Gracias a este mecanismo, el tamaño del espacio de direcciones de un proceso queda determinado por la arquitectura del sistema y la capacidad del almacenamiento secundario, en lugar de estar limitado únicamente por la cantidad de memoria principal disponible.

La memoria virtual introduce una abstracción entre la memoria utilizada por los procesos y la memoria física del sistema. Como consecuencia, es necesario distinguir una serie de conceptos relacionados con el direccionamiento de memoria que se utilizarán a lo largo del capítulo:

\begin{itemize}
  \item \textit{Espacio de direcciones:} conjunto de direcciones de memoria que un proceso puede utilizar durante su ejecución.
  \item \textit{Espacio de direcciones virtuales:} espacio de direcciones asignado a un proceso mediante el mecanismo de memoria virtual. Generalmente es mayor que la memoria principal disponible.
  \item \textit{Dirección virtual:} dirección perteneciente al espacio de direcciones virtuales, generada por el proceso y traducida automáticamente a una dirección física antes de acceder a la memoria.
  \item \textit{Dirección física (o real):} dirección correspondiente a una ubicación de la memoria principal (RAM), utilizada finalmente por el hardware para realizar el acceso.
\end{itemize}

\subsection{Hardware y estructuras de control}

Comparada con el particionado fijo y dinámico, la paginación y segmentación simple aportan dos avances importantes: 
\begin{enumerate}
  \item Traducción dinámica de direcciones: todas las referencias son lógicas y se traducen a físicas en tiempo de ejecución. Un proceso puede entrar y salir de memoria y ocupar regiones diferentes cada vez.
  \item Carga no contigua: un proceso se divide en piezas (páginas o segmentos) que no necesitan estar contiguas en memoria gracias a la tabla de páginas/segmentos.
\end{enumerate}

Lo interesante de esto es que no hace falta cargar todo el proceso. Si la pieza con la próxima instrucción y la que tiene el dato a acceder están en memoria principal, la ejecución puede continuar. Esto es exactamente lo que se busca con la memoria virtual.

Cuando un proceso se carga en memoria, el sistema operativo transfiere únicamente una o unas pocas piezas iniciales. El conjunto de piezas que permanece en memoria principal en un instante dado se denomina \emph{conjunto residente}.

Mientras todas las referencias caigan dentro del conjunto residente, todo va bien. Si el procesador encuentra una dirección lógica que no está en memoria, genera una interrupción por fallo de acceso a memoria (memory fault). El SO bloquea el proceso, emite una lectura de disco para traer la pieza faltante y despacha otro proceso mientras espera. Cuando llega por E/S, el proceso vuelve a estado Listo.

¿Tener operaciones E/S adicionales no es ineficiente al momento de ejecutar un proceso? Puede parecerlo, pero es eficiente si se gestiona bien.

El uso de memoria virtual tiene dos implicaciones fundamentales:
\begin{enumerate}
  \item Más procesos en memoria: al cargar solo parte de cada uno, caben más. Aumenta la probabilidad de que siempre haya un proceso en estado Listo y mejora el uso del procesador.
  \item Procesos más grandes que la memoria física: se levanta la restricción más básica de la programación. El programador ya no necesita hacer overlays manuales. Ve una memoria enorme (del tamaño del disco) y es el SO quien trae las piezas bajo demanda.
\end{enumerate}

Como un proceso se ejecuta únicamente si está en memoria principal, también se le llama memoria real. Pero la que percibe el programador, alojada en disco, es la \emph{memoria virtual}. Esto permite multiprogramación muy efectiva y libera al usuario de las limitaciones de la memoria física.

\subsection{Localidad y memoria virtual}

Los beneficios de la memoria virtual son claros, pero ¿es práctica? La experiencia demostró que sí, y hoy es un componente esencial de los SO modernos.

La clave está en esto: un proceso grande tiene mucho código y datos, pero en un intervalo corto la ejecución se concentra en una pequeña sección (ej: una subrutina) y en uno o dos arreglos. Sería un desperdicio cargar decenas de piezas cuando solo se usarán unas pocas.

Pero el SO debe ser inteligente. En estado estable, casi toda la memoria está ocupada, así que para traer una pieza debe expulsar otra. Si expulsa una que se necesitará enseguida, tendrá que traerla de vuelta al momento. El exceso de esto provoca \emph{thrashing} (hiperpaginación): el sistema pasa más tiempo intercambiando que ejecutando.

Evitar la hiperpaginación fue un gran tema de investigación en los 70 y llevó a algoritmos que, basados en la historia reciente, intentan adivinar qué piezas tienen menos probabilidad de usarse pronto.

Todo esto se basa en el principio de localidad: las referencias a programas y datos tienden a agruparse. Por eso es válido suponer que solo unas pocas piezas se necesitarán a corto plazo y se puede predecir cuáles.

Para que la memoria virtual sea efectiva se necesitan dos cosas:
\begin{enumerate}
  \item Soporte de hardware para paginación y/o segmentación.
  \item Software del SO para gestionar el movimiento de páginas/segmentos entre memoria secundaria y principal.
\end{enumerate}

\subsection{Paginación con memoria virtual}

El término memoria virtual suele asociarse a paginación. Con paginación simple, toda la memoria se divide en marcos del mismo tamaño y el proceso en páginas iguales, y se cargan todas las páginas en marcos no necesariamente contiguos. En paginación virtual, el principio es el mismo, pero no es necesario cargar todas las páginas para ejecutar.

Cada proceso sigue teniendo su propia tabla de páginas, pero las entradas (PTE - Page Table Entry) son más complejas añadiendo un bit de presencia, un bit de modificación y otros bits de control.

\subsubsection{Estructura de la tabla de páginas}

La tabla de páginas de un proceso es, conceptualmente, un vector o arreglo en el que el índice de la celda es el Número de Página Virtual ($p$).

Cada posición (llamada Page Table Entry o PTE) almacena:
\begin{itemize}
  \item Número de Marco Físico ($f$): El identificador del bloque real en memoria principal donde está guardada esa página.
  \item Bit de presencia: indica si la página está en memoria principal o secundaria.
  \item Bits de Control: Permisos de acceso (Lectura/Escritura/Ejecución), bit de modificación (dirty bit), entre otros.
\end{itemize}
Procedimiento de traducción en Paginación:
\begin{enumerate}
  \item La dirección virtual llega como el par ($p$,$d$), donde $p$ es la página y $d$ es el offset (desplazamiento).
  \item La MMU usa \(p\) para indexar la tabla del proceso activo y extrae el marco $f$.
  \item Como las páginas y los marcos tienen exactamente el mismo tamaño, el desplazamiento dentro de la página es idéntico al desplazamiento dentro del marco.
  \item La dirección física resultante es simplemente la combinación del marco con el offset: ($f$,$d$).
\end{enumerate}

Como la tabla es de longitud variable, debe estar en memoria principal, no en registros. Un registro guarda la dirección de inicio de la tabla del proceso en ejecución.

Normalmente hay una tabla por proceso, pero en sí mismas, las tablas pueden ser enormes. Por ejemplo: con 2GB de memoria virtual y páginas de 512 bytes, se necesitan \(2^{22}\) entradas. Ocuparía demasiada memoria. La solución es guardar las propias tablas de páginas en memoria virtual, por lo que también se paginan. Ahora el problema reside en cómo encontrar la tabla de un proceso si para hallarla debería consultarse la memoria virtual, y por ende la tabla de páginas.

Para solucionar el problema anterior y manejar tablas grandes se usa el esquema de dos niveles: un directorio de páginas permanente en memoria real. Cada entrada del directorio apunta a una tabla de páginas. Si el directorio tiene \(X\) entradas y cada tabla \(Y\) entradas, el proceso puede tener hasta \(X \times Y\) páginas. Normalmente cada tabla ocupa como máximo una página.

\begin{figure}[!ht]
  \centering
  \includegraphics[width=0.9\textwidth]{chapters/05_unit/media/dirs_virtual_mem.pdf}
  \caption{Esquema de dos niveles.}
  \label{fig:dirs_virtual_mem}
\end{figure}

Ejemplo de 32 bits con páginas de 4KB (\(2^{12}\)) (Figura~\ref{fig:dirs_virtual_mem}): el espacio virtual de 4GB (\(2^{32}\)) tiene \(2^{20}\) páginas. Si cada PTE ocupa 4 bytes, la tabla de usuario tendría \(2^{20}\) entradas = 4MB (\(2^{10}\) páginas). Esa tabla de 4MB puede mantenerse en memoria virtual y ser mapeada por una tabla raíz de \(2^{10}\) entradas = 4KB, que sí permanece siempre en memoria real.

Traducción en dos niveles (Figura~\ref{fig:addressing_virtual_mem}): los primeros 10 bits de la dirección virtual indexan la tabla raíz para encontrar la página de la tabla de usuario. Si esa página no está en memoria, hay fallo de página. Si está, los siguientes 10 bits indexan la tabla de usuario para encontrar el marco de la página referenciada.

\begin{figure}[!ht]
  \centering
  \includegraphics[width=0.9\textwidth]{chapters/05_unit/media/addressing_virtual_mem.pdf}
  \caption{Traducción de dirección en un sistema de paginación con directorios.}
  \label{fig:addressing_virtual_mem}
\end{figure}

\subsubsection{Buffer de traducción anticipada (TLB)}

En principio, cada referencia virtual requiere dos accesos físicos: uno para buscar la entrada en la tabla de páginas y otro para el dato. Esto duplicaría el tiempo de acceso.

Para evitarlo, se usa una caché de alta velocidad para entradas de tabla: el TLB (Translation Lookaside Buffer). Guarda las entradas usadas más recientemente.

El funcionamiento es así:
\begin{enumerate}
  \item Dada una dirección virtual, el procesador consulta el TLB. Si hay acierto (hit), obtiene el número de marco y forma la dirección física.
  \item Si hay fallo (miss), busca la entrada en la tabla de páginas en memoria. Si el bit de presencia está activo, la página está en memoria: toma el marco, forma la dirección real y actualiza el TLB con la nueva entrada.
  \item Si el bit de presencia no está activo, ocurre un fallo de página (page fault) y el SO debe traer la página del disco.
\end{enumerate}

Por el principio de localidad, la mayoría de las referencias usan páginas recientes, por lo que la mayoría son aciertos en el TLB, mejorando mucho el rendimiento.

Detalles de implementación: como el TLB solo tiene parte de la tabla, no se puede indexar solo por número de página. Cada entrada del TLB debe incluir el número de página completo. El hardware compara en paralelo todas las entradas (mapeo asociativo), a diferencia del mapeo directo de la tabla de páginas.

Interacción con la caché: la dirección virtual se traduce primero vía TLB/tabla a dirección física. Luego, con esa dirección física, se consulta la caché de memoria principal para ver si el bloque ya está allí. Por lo tanto, una sola referencia puede implicar buscar en TLB, en tabla de páginas, en caché, en memoria principal y hasta en disco.

\subsubsection{Tamaño de página}

El tamaño de página es una decisión clave de hardware con varios factores en conflicto.

El compromiso según tamaño de página puede categorizarse de la siguiente manera:
\begin{itemize}
  \item Página pequeña: menos fragmentación interna y aprovecha mejor la memoria principal.
  \item Página grande: menos páginas por proceso, por lo tanto tablas de páginas más pequeñas. Si las páginas son muy pequeñas, las tablas se hacen tan grandes que parte de ellas debe quedar en memoria virtual, pudiendo causar doble fallo de página: uno para traer la porción de la tabla y otro para traer la página del proceso. Además, los dispositivos de almacenamiento rotacionales transfieren bloques grandes de forma más eficiente, lo que favorece páginas grandes.
\end{itemize}

Por el principio de localidad, al aumentar moderadamente el tamaño de página, cada fallo incorpora más datos cercanos a la referencia realizada, reduciendo la tasa de fallos. Sin embargo, si las páginas son demasiado grandes, con una cantidad fija de memoria principal cabrán menos páginas en el conjunto residente, aumentando los reemplazos y, en consecuencia, la tasa de fallos. Por ello, existe un tamaño de página que representa un compromiso entre ambos efectos.

Además, la tasa de fallos también depende de la política de asignación de memoria: cuanto mayor sea el número de marcos asignados a un proceso, menor será, en general, la probabilidad de producir un fallo de página.

Aunque la memoria principal crece, el tamaño de los programas también y la localidad disminuye por técnicas modernas:
\begin{itemize}
  \item Programación orientada a objetos: muchas referencias dispersas en muchos objetos pequeños.
  \item Multithreading: cambios abruptos en el flujo de instrucciones.
\end{itemize}

Con un TLB de tamaño fijo, si los procesos crecen y la localidad baja, la tasa de aciertos del TLB cae y se vuelve un cuello de botella. Agrandar el TLB tiene límites por su interacción con la caché y el hardware. Usar páginas más grandes haría que cada entrada del TLB cubra más memoria, pero como vimos, degrada el rendimiento.

La solución consiste en utilizar múltiples tamaños de página.
Varias arquitecturas soportan tamaños múltiples. Esto permite más permite flexibilidad; regiones grandes y contiguas como el código se mapean con pocas páginas grandes, y la pila de hilos con páginas pequeñas. Esto mejora el rendimiento del TLB.

Sin embargo, la mayoría de los SO comerciales aún soportan un solo tamaño, porque cambiar a múltiples tamaños implica modificar muchos aspectos del sistema operativo y es muy complejo.

\subsection{Segmentación con memoria virtual}

La segmentación permite al programador ver la memoria como múltiples espacios de direcciones o segmentos de tamaño variable y dinámico. Las referencias son del tipo \texttt{(número de segmento, offset)}.

Ventajas para el programador:
\begin{enumerate}
  \item Estructuras de datos que crecen: no hace falta adivinar el tamaño. Se le asigna su propio segmento y el SO lo agranda o achica según necesidad. Si no hay espacio en memoria principal, el SO lo mueve a un área más grande o lo intercambia a disco.
  \item Recompilación independiente: los programas pueden modificarse y recompilarse por separado sin tener que re-enlazar y recargar todo.
  \item Compartición: un programa utilitario o tabla de datos puede ponerse en un segmento compartido por varios procesos.
  \item Protección: al tener un conjunto bien definido de programas o datos por segmento, es más fácil asignar privilegios de acceso.
\end{enumerate}

\subsubsection{Organización}

Como en segmentación simple, cada proceso tiene su tabla de segmentos con dirección de inicio y longitud de cada segmento. Con memoria virtual, las entradas son más complejas, porque no todos los segmentos están en memoria:
\begin{itemize}
  \item Bit de presencia: indica si el segmento está en memoria principal. Si está, la entrada guarda su dirección de inicio y longitud.
  \item Bit de modificación: indica si fue alterado desde que se cargó. Si no cambió, no hace falta escribirlo a disco al reemplazarlo.
  \item Otros bits de control (protección, compartición).
\end{itemize}

Traducción de direcciones: la dirección lógica \texttt{(segmento, offset)} se traduce a física usando la tabla de segmentos. Como la tabla es de longitud variable, está en memoria principal y un registro guarda su dirección de inicio cuando el proceso corre.

El número de segmento indexa la tabla para obtener la dirección de inicio del segmento, se verifica que el offset sea menor que la longitud, y se suma el offset para obtener la dirección real:
\begin{center}
\texttt{Dirección física = Dirección de inicio del segmento + offset}
\end{center}

\subsection{Paginación y Segmentación Combinadas}

Paginación y segmentación tienen fortalezas distintas. La paginación es transparente al programador, elimina la fragmentación externa y permite algoritmos de gestión sofisticados al mover piezas fijas. La segmentación es visible al programador y aporta manejo de estructuras que crecen, modularidad, compartición y protección. Para unir ambas ventajas, algunos sistemas combinan las dos.

\subparagraph{Funcionamiento} el espacio del usuario se divide en segmentos (a criterio del programador). Cada segmento, a su vez, se divide en páginas de tamaño fijo, igual al marco de memoria. Si un segmento es menor que una página, ocupa solo una. Para el programador la dirección lógica sigue siendo \texttt{número de segmento + offset}, pero el sistema interpreta ese offset como \texttt{número de página + offset dentro del segmento}.

\begin{figure}[!ht]
  \centering
  \includegraphics[width=0.9\textwidth]{chapters/05_unit/media/comb_page_segment_vm.pdf}
  \caption{Traducción de direcciones en un sistema con paginación y segmentación.}
\end{figure}

\subparagraph{Traducción} cada proceso tiene una tabla de segmentos y una tabla de páginas por cada segmento. Un registro guarda la dirección de la tabla de segmentos del proceso en ejecución. El procesador usa el número de segmento para indexar la tabla de segmentos y hallar la tabla de páginas de ese segmento. Luego usa el número de página para indexar esa tabla y obtener el marco. Finalmente combina \texttt{marco + offset} para la dirección física.

\subparagraph{Formato de entradas} la entrada de la tabla de segmentos contiene la longitud del segmento y la base, que ahora apunta a una tabla de páginas. No necesita bits de presente/modificado, porque eso se maneja a nivel de página. La entrada de tabla de páginas es como en paginación pura: número de marco si está presente, bit de modificado, y bits de protección.

\subsubsection{Protección y Compartición}

La segmentación facilita la protección y compartición. Como cada entrada de segmento incluye longitud y base, un programa no puede acceder fuera de los límites del segmento. Para compartir, un mismo segmento puede ser referenciado en las tablas de segmentos de varios procesos. Esto también es posible con paginación, pero al no ser visible la estructura para el programador, es más difícil de especificar.

Un esquema común es la estructura de anillos: los anillos interiores (número más bajo) tienen más privilegios. Típicamente el anillo 0 es para el kernel del SO y los anillos exteriores para aplicaciones.

Principios básicos:
\begin{itemize}
  \item Un programa solo puede acceder a datos del mismo anillo o de un anillo menos privilegiado.
  \item Un programa puede llamar a servicios del mismo anillo o de un anillo más privilegiado.
\end{itemize}

